% !TeX root = ../../../../../main.tex
% chapters/sections/chapter5/subsections/a_shock/subsubsections/interpretations/rise.tex

\subsubsection{\(i_t^H\)が上昇する政策群（TR-PPI, POLT, IT-PPI, PPLT）の解釈}
\label{sec:chapter5_a_shock_interpretations_rise}
\begin{enumerate}
    \item
    \label{itm:chapter5_a_shock_interpretations_rise_first}
        \textbf{1期：\(y_t^H,c_t^{H\to W}\)の下落}
        \begin{description}[style=nextline]
            \item[\(\hat{a}_1^H\)の負への下落と\(i_1^H\)の上昇]
            \label{itm:chapter5_a_shock_interpretations_rise_first_a_H}
                \(a\)ショックが発生し、予期せず\(\hat{a}_1^H\)が負へ大きく下落する。
                負へ下落した\(\hat{a}_t^H\)は自己回帰過程の式
                \eqref{form:chapter4_shocks_a_eq}
                にしたがって増加し、やがて初期定常状態\(0\)へ回帰する。
                また後述の諸変数の変動により\(i_1^H\)は金融政策の式を通じて上昇する。
                上昇した\(i_t^H\)は金融政策の式にしたがって将来のいずれかの期から低下しはじめ、
                やがて初期定常状態\(i_{ss}^H\)へ回帰する。
            \item[\(\hat{a}_1^H\)が負になることによるYS-\(\hat{p}\),CS-\(\hat{p}\)曲線の上移動]
            \label{itm:chapter5_a_shock_interpretations_rise_first_a_H_ys_cs}
                \(\hat{a}_1^H\)が負になることを受けて供給側では以下のような反応が起こる。
                \(\hat{p}_0^H=0\)であり、また\(\hat{a}_1^H\)が動き始めるこの時点においてまだ
                \(\operatorname{E}_1[\hat{\pi}_2^H]=0\)である。
                したがってYS-\(\hat{p}\),CS-\(\hat{p}\)曲線の切片の主決定要因
                \(\hat{p}_0^H+\beta_{ss}^H\operatorname{E}_1[\hat{\pi}_2^H]-2\kappa^H\hat{a}_1^H
                =-2\kappa^H\hat{a}_1^H\)
                が正となる。
                よってYS-\(\hat{p}\),CS-\(\hat{p}\)曲線は切片が上昇し、上へ移動する。
                これにより\(y_1^H,c_1^{H\to W}\)は下落し、\(p_1^H,p_1^{H\to W}\)は上昇する。
                また完全伸縮物価\(\xi^H=0\)を仮定した場合、
                YS-\(\hat{p}\),CS-\(\hat{p}\)曲線は垂直となり
                上記の\(\hat{p}\)軸切片の上昇は\(y_1^H\)軸,\(c_1^{H\to W}\)軸切片の下落となる。
                よってYS-\(\hat{p}\),CS-\(\hat{p}\)曲線は左移動し、
                \(y_1^{H,pot}\)は\(y_1^H\)よりも大きく下落する。
            \item[\(\hat{p}_1^H\)が正になることによるYS-\(\hat{p}\),CS-\(\hat{p}\)曲線のさらなる上移動]
            \label{itm:chapter5_a_shock_interpretations_rise_first_p_H_YS_CS}
                \(\hat{\pi}_1^H=\hat{p}_1^H-\hat{p}_0^H=\hat{p}_1^H\)が正となった。
                そこで\(\operatorname{E}_1[\hat{\pi}_2^H]\)も正であると仮定する。
                するとYS-\(\hat{p}\),CS-\(\hat{p}\)曲線の切片の主決定要因
                \(\hat{p}_0^H+\beta_{ss}^H\operatorname{E}_1[\hat{\pi}_2^H]-2\kappa^H\hat{a}_1^H
                =\beta_{ss}^H\operatorname{E}_1[\hat{\pi}_2^H]-2\kappa^H\hat{a}_1^H
                \approx\operatorname{E}_1[\hat{\pi}_2^H]-2\kappa^H\hat{a}_1^H\)
                がさらに上昇する。
                よってYS-\(\hat{p}\),CS-\(\hat{p}\)曲線は切片がさらに上昇し、さらに上へ移動する。
                これにより\(y_1^H,c_1^{H\to W}\)はさらに減少し、\(p_1^H,p_1^{H\to W}\)はさらに増加する。
                また
                \ref{itm:chapter5_a_shock_interpretations_rise_first_a_H_ys_cs}
                の末尾と同様の理由により\(y_1^{H,pot}\)は\(y_1^H\)よりも大きく下落する。
            \item[\(i_1^H\)が上昇したことによるYD,CD曲線の下移動]
            \label{itm:chapter5_a_shock_interpretations_rise_first_i_H_YD_CD}
                このとき需要側においても変化が生じる。
                \(p_1^H,p_1^{H\to W}\)の上昇、
                あるいは\(y_t^{H,pot}\)が\(y_1^H\)よりも大きく下落したことにより、
                \(i_1^H\)は金融政策の式を通じて上昇する。
                するとYD,CD曲線の係数の主決定要因
                \(\frac{1}{(1+i_1^H)\operatorname{E}_1[\lambda_2^H]}\)
                が下落する。
                よってYD,CD曲線は係数が下落し、下へ移動する。
                これにより\(y_1^H,c_1^{H\to W}\)および\(p_1^H,p_1^{H\to W}\)は下落する。
            \item[\(\operatorname{E}_1[\lambda_2^H]\)の上昇によるYD,CD曲線のさらなる下移動]
            \label{itm:chapter5_a_shock_interpretations_rise_first_E_lambda_YD_CD}
                自国コンティンジェント債券に関する前方展開されたオイラー方程式
                \eqref{form:chapter5_beta_shock_policies_overview_euler_contingent_lim}
                により、\(\operatorname{E}_1[\lambda_2^H]\)は
                数列(\(i_{1+k}^H\))と\(\lim_{k\to\infty}p_{1+k+1}^{H\to W}\)によって一意に決定されるが、
                これらは金融政策の式により一意に定められる。
                現実において数列(\(i_{1+k}^H\))については
                政策によりある程度の予想ができると考えられることから、
                ここでは数列(\(i_{1+k}^H\))は正確に知っているものとする。
                一方で\(\lim_{k\to\infty}p_{1+k+1}^{H\to W}\)を正確に知ることは困難であるため、
                ここでは
                \(\lim_{k\to\infty}p_{1+k+1}^{H\to W}\)を
                \(p_{ss}^{H\to W}\)として予想するものとする。
                すると
                \(\lim_{k\to\infty}p_{1+k+1}^{H\to W}\)は
                定常状態を維持する。
                また数列(\(i_{1+k}^H\))全体の上昇は右辺の無限乗積
                \(\operatorname{E}_1\left[\prod_{k=1}^{\infty}\beta_{ss}^H(1+i_{1+k}^H)\right]\)
                を増加させる。
                したがって右辺は定常状態から上昇するのだから
                \(\operatorname{E}_1[\lambda_2^H]\)
                は上昇する。
                するとYD,CD曲線の係数の主決定要因
                \(\frac{1}{(1+i_1^H)\operatorname{E}_1[\lambda_2^H]}\)
                はさらに下落する。
                よってYD,CD曲線は係数がさらに下落し、さらに下へ移動する。
                これにより\(y_1^H,c_1^{H\to W}\)および\(p_1^H,p_1^{H\to W}\)はさらに下落する。
            \item[1期のまとめ]
            \label{itm:chapter5_a_shock_interpretations_rise_first_summary}
                \(\hat{a}_1^H\)が負になることによるYS-\(\hat{p}\),CS-\(\hat{p}\)曲線の上移動
                \(\hat{p}_1^H\)が正になることによるYS-\(\hat{p}\),CS-\(\hat{p}\)曲線のさらなる上移動
                \(i_1^H\)が上昇したことによるYD,CD曲線の下移動
                \(\operatorname{E}_1[\lambda_2^H]\)の上昇によるYD,CD曲線のさらなる下移動
                これらを合わせることにより
                YS-\(\hat{p}\),CS-\(\hat{p}\)曲線は大きく上移動
                YD,CD曲線は大きく下移動
                となる。その結果として
                \(y_1^H,c_1^{H\to W}\)は下落する。
                一方で価格については、
                供給曲線の上移動による上昇圧力と
                需要曲線の下移動による下落圧力が拮抗する。
                CS-\(\hat{p}\)曲線はYS-\(\hat{p}\)曲線よりも傾きが大きいため、
                需要曲線の上下移動による圧力をより強く受ける。
                そのため\(p_t^{H\to W}\)については
                CS-\(\hat{p}\)曲線の上移動よりも需要曲線の下移動の圧力を強く受け下落する。
                一方で\(p_t^H\)については
                YS-\(\hat{p}\)曲線の上移動と需要曲線の下移動の圧力が拮抗し、
                POLTやTR-PPI、IT-PPIにおいてはYS-\(\hat{p}\)曲線の上移動による圧力がやや勝り上昇する。
                しかしPPLTは\(p_t^H\)を直接の目標とするため、
                \(p_t^H\)を維持するように過不足なく数列(\(i_{t+k}^H\))の全体を引き上げる。
                これにより需要曲線がPOLTやTR-PPI、IT-PPIよりもさらに下移動し\(p_t^H\)はほぼ維持される。
        \end{description}
\textbf{2期以降：\(y_t^H,c_t^{H\to W}\)の単調増加による初期定常状態への収束}
        \begin{description}[style=nextline]
            \item[負の\(\hat{a}_t^H\)の増加と\(i_t^H\)の減少]
            \label{itm:chapter5_a_shock_interpretations_rise_second_a_H_i_H}
                負の\(\hat{a}_t^H\)は自己回帰過程の式
                \eqref{form:chapter4_shocks_a_eq}
                にしたがい\(0\)へ向かって増加する。
                \(i_t^H\)は減少していく。
            \item[\(\hat{y}_t^H-\hat{a}_t^H\)が正で\(0\)に収束することにより\(\operatorname{E}_t[\hat{\pi}_{t+1}^H]\)が減少幅を縮小していくこと]
            \label{itm:chapter5_a_shock_interpretations_rise_second_y_H_YS_CS}
                供給側では以下のような反応が起こる。
                \(\hat{a}_t^H\)と\(\hat{y}_t^H\)はともに負になった。
                ここで\(\hat{y}_t^H-\hat{a}_t^H\)は正とする。
                またDynareのシミュレーションの仕組みにより
                \(\hat{y}_t^H\)は\(0\)に収束するから\(\hat{y}_t^H-\hat{a}_t^H\)も\(0\)に収束する。
                よってYS-\(\pi\)曲線（差分形式）において右辺の主決定要因
                \(\hat{y}_t^H-\hat{a}_t^H\)
                が正で\(0\)に収束するから、左辺の
                \(\hat{\pi}_t^H-\beta_{ss}^H\operatorname{E}_t[\hat{\pi}_{t+1}^H]
                =\operatorname{E}_{t-1}[\hat{\pi}_t^H]-\beta_{ss}^H\operatorname{E}_t[\hat{\pi}_{t+1}^H]
                \approx\operatorname{E}_{t-1}[\hat{\pi}_t^H]-\operatorname{E}_t[\hat{\pi}_{t+1}^H]\)
                も正で\(0\)に収束し、
                よって\(\operatorname{E}_t[\hat{\pi}_{t+1}^H]\)は減少幅を縮小していく。
            \item[\(\operatorname{E}_t[\hat{\pi}_{t+1}^H]\)が減少幅を縮小していくことと\(\hat{a}_t^H\)が増加幅を縮小していくことによるYS-\(\hat{p}\),CS-\(\hat{p}\)曲線の上移動の縮小]
            \label{itm:chapter5_a_shock_interpretations_rise_second_p_H_YS_CS}
                \(\operatorname{E}_t[\hat{\pi}_{t+1}^H]
                =\hat{\pi}_{t+1}^H\)
                が減少していくということは\(\hat{p}_{t+1}^H\)の増加幅が縮小していくことを意味する。
                また\(\hat{a}_t^H\)は増加幅を縮小していくのだから
                正の\(-2\kappa^H\hat{a}_t^H\)も減少幅を縮小していく。
                するとYS-\(\hat{p}\),CS-\(\hat{p}\)曲線の切片の主決定要因
                \(\hat{p}_{t-1}^H+\beta_{ss}^H\operatorname{E}_t[\hat{\pi}_{t+1}^H]-2\kappa^H\hat{a}_t^H\)
                において
                \(\hat{p}_{t-1}^H\)が増加幅を縮小していき、
                \(\beta_{ss}^H\operatorname{E}_t[\hat{\pi}_{t+1}^H]\)が減少幅を縮小していき、
                \(-2\kappa^H\hat{a}_t^H\)が減少幅を縮小していく。
                よって主決定要因は増加幅を縮小していくか、あるいは減少幅を縮小していく。
                主決定要因が増加幅を縮小していく場合、
                YS-\(\hat{p}\),CS-\(\hat{p}\)曲線は切片が増加幅を縮小していき、上移動を縮小していく。
                よって
                \(y_t^H,c_t^{H\to W}\)は減少幅を縮小していき、
                \(p_t^H,p_t^{H\to W}\)は増加幅を縮小していく。
                主決定要因が減少幅を縮小していく場合、
                YS-\(\hat{p}\),CS-\(\hat{p}\)曲線は切片が減少幅を縮小していき、下移動を縮小していく。
                よって
                \(y_t^H,c_t^{H\to W}\)は増加幅を縮小していき、
                \(p_t^H,p_t^{H\to W}\)は減少幅を縮小していく。
            \item[\(i_t^H\)が減少幅を縮小していくことによりYD,CD曲線が上移動を縮小していくこと]
            \label{itm:chapter5_a_shock_interpretations_rise_second_E_lambda_YD_CD}
                \(\beta_{ss}^H(1+i_{ss}^H)=1\)であり
                \(i_t^H\)は上から\(i_{ss}^H\)に向かって減少幅を縮小していくのだから、
                \(\beta_{ss}^H(1+i_{t+1}^H)\)は上から\(1\)に向かって減少幅を縮小していく。
                したがって自国コンティンジェント債券に関するオイラー方程式
                \eqref{form:chapter5_beta_shock_policies_overview_euler_contingent}
                より\(\lambda_t^H\)は減少幅を縮小していき、よってYD,CD曲線の係数の主決定要因
                \(\frac{1}{\lambda_t^H}\)
                は増加幅を縮小していく。
                よってYD,CD曲線は係数が増加幅を縮小していき、上移動を縮小していく。
                これにより\(y_t^H,c_t^{H\to W}\)および\(p_t^H,p_t^{H\to W}\)は増加幅を縮小していく。
            \item[2期以降のまとめ]
            \label{itm:chapter5_a_shock_interpretations_rise_second_summary}
                \(\operatorname{E}_t[\hat{\pi}_{t+1}^H]\)が減少幅を縮小していくことと
                \(\hat{a}_t^H\)が増加幅を縮小していくことにより
                YS-\(\hat{p}\),CS-\(\hat{p}\)曲線は上移動を縮小していくか下移動を縮小していく。
                \(i_t^H\)が減少幅を縮小していくことによりYD,CD曲線は上移動を縮小していく。
                これらにより\(y_t^H,c_t^{H\to W}\)は増加幅を縮小していく。
                一方で価格については、
                供給曲線の下移動による下落圧力と
                需要曲線の上移動による上昇圧力が拮抗する。
                1期と同様に、CS-\(\hat{p}\)曲線はYS-\(\hat{p}\)曲線よりも傾きが大きいため、
                需要曲線の上下移動による圧力をより強く受ける。
                そのため\(p_t^{H\to W}\)については
                CS-\(\hat{p}\)曲線の下移動よりも需要曲線の上移動による圧力を強く受け、
                需要曲線が上移動を縮小していくにともない増加幅を縮小していく。
                一方で\(p_t^H\)については
                CS-\(\hat{p}\)曲線の下移動と需要曲線の上移動の圧力が拮抗し、
                政策や時期により異なる動きを見せる。
        \end{description}
\end{enumerate}